\section{Results}
\subsection{Clinical framework and analysis populations}
The source framework combined at least one physician criterion with one patient criterion at W4, W12 and W24; W2 assessed itch alone (Fig.~\ref{fig:Figure1}). We used these stage-specific rules to calculate recorded-score targets. The serum-assisted EASI readout was evaluated as a parallel research measure and did not determine dose-reduction decisions.

% INSERT: Figure1


The clinical comparison included 163 patients: 65 started dupilumab, 56 abrocitinib 100 mg and 42 abrocitinib 200 mg. Median baseline EASI was 15.20, 11.70 and 19.35, respectively; age, sex and BMI also differed (Table~\ref{tab:baseline}). Clinical treatment comparisons and linked-serum modelling used separate eligibility rules. Their populations and overlap are shown in Supplementary Fig.~\suppfig{39}.

% INSERT: Table1

\subsection{Clinical responses differed across initial treatment groups}
Clinical activity decreased and patients reached the successive treatment targets (Fig.~\ref{fig:Figure2}). At week 2, PP-NRS was at most 4 in 15/27 evaluable dupilumab recipients, 34/37 receiving initial abrocitinib 100 mg and 33/36 receiving initial abrocitinib 200 mg. The adjusted 100 mg--dupilumab difference in this early itch-target rate was 41.7 percentage points (95\% CI, 15.1--69.4; Holm-adjusted $P=0.0095$); adding baseline itch gave 41.4 points (13.0--66.8). W2 itch was unavailable for 38/65, 19/56 and 6/42 patients, respectively, so these available-case associations remain sensitive to follow-up selection. Weighted worst-case bounds at W2 were 19.4--82.0\% for dupilumab and 76.0--94.8\% for A100, allowing the contrast to reverse under extreme assignments of missing outcomes. These are sensitivity bounds, not confidence intervals. At week 24, the combined physician- and patient-reported target was reached by 26/36 (72.2\%), 22/24 (91.7\%) and 20/23 (87.0\%), respectively. Allowing all unknown outcomes to fail or succeed gave W24 combined-target bounds of 40.0--84.6\%, 39.3--96.4\% and 47.6--92.9\%. All specified clinical trajectories, component targets and exact denominators are retained in Supplementary Figs.~\suppfig{1}--\suppfig{4}, \suppfig{15} and \suppfig{23}--\suppfig{26}.

% INSERT: Figure2

The lower observed EASI in the 100 mg group did not persist as a clear dose-group difference after baseline adjustment. The adjusted 100-minus-200 mg difference was 0.62 points at week 12 (95\% CI, $-2.08$ to 3.33) and 0.34 at week 24 ($-1.10$ to 1.77). Weighting improved full-cohort baseline balance, although residual imbalance remained in endpoint-specific available-case samples (Supplementary Fig.~\suppfig{40}). These comparisons describe initial treatment under subsequent clinical management, including later regimen changes. Three W0 regimen entries conflicted with the initial-dose label and were retained as unresolved source conflicts. Recorded exposure composition and post-response reduction pathways are shown separately (Supplementary Figs.~\suppfig{16}--\suppfig{17} and \suppfig{41}). Safety-related documentation through W24 was available for 27/65, 20/56 and 19/42 patients, respectively (Supplementary Table~S3). These counts describe note availability; event incidence, severity and event-time treatment attribution were not adjudicated.

\subsection{Molecular trajectories differed across treatment courses}
All 23 measured serum markers were retained, including CCL18, IL-13, IL-22 and IL-31. The W0/W12/W24 profiles showed treatment-dependent molecular changes alongside clinical improvement (Fig.~\ref{fig:Figure3}). Fixed three-visit displays used the same patients within each marker and treatment; the complete window analyses used all available pairs. Thus, differences in their sample sizes reflect distinct eligible sets, rather than changes to the measurements.

% INSERT: Figure3

For CCL17 over W0--W12, the adjusted dupilumab-minus-A100 difference was $-1.07$ log$_2$ units (95\% CI, $-1.67$ to $-0.48$; q=0.0109). The assay-batch sensitivity estimate was $-0.72$, with an interval crossing zero (q=0.369), so this contrast was not robust to that specification. Direct adjusted contrasts, rather than separate within-group P values, addressed treatment differences. All windows and complete comparisons are retained in Supplementary Figs.~\suppfig{5}--\suppfig{14}.

Molecular and clinical changes were associated for selected marker--outcome combinations. In A100 over W0--W12, the IL-5 and PP-NRS changes had Spearman $\rho=0.742$ (31 pairs; panel-wide q=0.00156). The complete declared association family also identified TNF-$\alpha$ and TSLP changes with EASI change over W0--W24. Their descriptive response-category profiles are shown in Supplementary Fig.~\suppfig{43}, alongside the full-panel response-depth results in Supplementary Figs.~\suppfig{27}--\suppfig{34}. These categories are ordered clinical-improvement groups, not time points or independent validation samples.

The fixed-patient IgE check used the same 11/9/10 patients across W0/W12/W24 (Supplementary Fig.~\suppfig{21}). It showed decreasing group medians despite differences in the visit-level summaries, illustrating sensitivity to sample composition without proving its cause. Dupilumab W4 eosinophil observations were absent; no W4 value was interpolated. These analyses distinguish a broad molecular response from the narrower task of estimating an EASI value with three markers.

\subsection{The common-model procedure estimated follow-up EASI closely}
The model dataset contained 316 complete-input visits from 160 patients: 152 at baseline, 89 at week 12 and 75 at week 24. Development used 241 baseline/week-12 visits from 156 patients. The final equation combined BSA with CCL17, CCL26 and IL-5 through the same calculation before and during treatment. It required no treatment label, visit week or previous score (Fig.~\ref{fig:Figure4}).

Under nested patient-level validation, MAE was 3.19 at baseline (95\% CI, 2.75--3.70), 1.30 at week 12 (0.96--1.71) and 1.19 at week 24 (0.89--1.53). Spearman correlations were 0.918, 0.927 and 0.904, and concordance correlations were 0.910, 0.900 and 0.869 (Table~\ref{tab:model}). Mean signed errors at W12/W24 were +0.82/+0.84 EASI points, and calibration slopes were 0.81 (95\% CI, 0.66--0.94) and 0.73 (0.59--0.82), indicating residual overestimation and scale miscalibration. The cumulative error distributions show that 41.8\%, 84.3\% and 82.0\% of predictions, respectively, were within two EASI points. These are proportions of individual predictions, not agreement between group means.

% INSERT: Table2
% INSERT: Figure4

Baseline severe disease remained the principal accuracy boundary: MAE was 5.38 among the 44 patients with EASI $>21$, compared with 1.15 and 2.55 in the two lower strata (Supplementary Fig.~\suppfig{37}). There was only one such severe observation at week 12 and none at week 24. Centre and treatment holdouts also showed variation (Supplementary Fig.~\suppfig{22}). Matched BSA-area control MAEs were 1.18 and 1.02 at W12/W24, compared with 1.30 and 1.19 for the primary procedure. Paired excess errors were 0.12 points (95\% CI, 0.04--0.20) and 0.17 (0.11--0.24), conditional on saved predictions. Thus, adding the serum inputs did not improve current EASI estimation in this evaluation. Complete serum-only and fixed-structure controls remain in the source tables.

Removing eight W0 and two W12 reference-inconsistent visits from training, while retaining the original folds and evaluation outcomes, gave sensitivity MAEs of 3.19/1.15/1.06 at W0/W12/W24. This refit improved follow-up error but did not establish molecular added accuracy over the identically refitted BSA-area control. Follow-up calibration slopes remained below one. The frozen primary equation and primary estimates were retained (Supplementary Table~S4).

Independent models for PP-NRS, adult DLQI and recorded ordinal IGA did not reproduce the EASI performance with this restricted input panel. Week-12/week-24 MAEs were 2.78/2.53 for PP-NRS and 5.90/6.28 for DLQI; exact IGA agreement was 35.4\%/28.0\% (Supplementary Fig.~\suppfig{35}).

\subsection{The equation tracked within-person change with window-dependent precision}
Estimated and observed EASI changes followed similar ordering across all three paired intervals (Fig.~\ref{fig:Figure5}). Change MAEs were 3.87 for baseline--week 12 (85 patients; 95\% CI, 3.08--4.70), 3.47 for baseline--week 24 (67 patients; 2.68--4.39) and 1.49 for weeks 12--24 (55 patients; 1.00--2.19). The corresponding Spearman correlations were 0.873, 0.940 and 0.900. The two intervals beginning at baseline retained larger errors despite their high correlations; estimation of the baseline score contributed to this difference.

% INSERT: Figure5

Using estimated scores at both visits to identify EASI75 gave sensitivities of 76.5\% at week 12 (95\% CI, 65.6--87.0) and 86.1\% at week 24 (76.9--94.0), with specificities of 98.6\% (95.2--100) and 96.2\% (87.5--100). Only 36 and 13 observed nonresponders informed specificity at the two visits. Some observed responders were missed despite the high specificity. Sensitivity for the deeper EASI90 response was lower, at 63.0\% and 52.9\% (Supplementary Fig.~\suppfig{38}). These separate paired evaluations support assessment of change alongside the absolute reading; they do not establish a universally interchangeable measure of treatment response.
